\subsection{One-step supervision and independently sampled positive banks}
The original shuffled training design placed the scheduled structure in the first positive slot. Our paired arm kept the identical sequence of composition labels, drew all four positive indices independently with replacement in each block, and generated independent query and negative sets. Hashes of the complete ordered schedules agree for each seed and budget. The distinction matters when a composition has more than one observed geometry. The independent-positive design satisfies the bank-exchangeability hypothesis of the supplement conditional on the labels, parameters and query inputs.
We trained a dedicated one-step control using the MeanFlow identity \cite{meanflow} on the same set network. Two independent uniform times are sorted into $r\leq t$, with probability $1/4$ assigned to the diagonal $r=t$. On the straight log-lattice path from data at $0$ to noise at $1$, let $v_k=k_N-k_D$. On the site torus we shuffle target sites within species, select componentwise wrapped noise--data displacements, subtract their occupied-site mean, and call the resulting tangent $v_f$. The network predicts $u=(u_k,u_f)$, with the site head expressed in units of $\sqrt{12}$. Its target is
\[ \operatorname{sg}\!\left[(v_k,\sqrt{12}\,v_f)-(t-r)\{\partial_k u\,v_k+\partial_f u\,v_f+\partial_t u\}\right]. \]
The derivative follows $k_t=k_N+(t-1)v_k$ and $f_t=f_N+(t-1)v_f\pmod 1$ with $r$ fixed and is evaluated by a forward-mode Jacobian--vector product. Centering $v_f$ translates all occupied clean sites by the same vector at $t=0$, preserving their relative positions. Six interval coefficients feed $(t-r)$ into the shared lattice input features. Training uses the same lattice and masked site squared errors as compact flow. At inference one call at $(r,t)=(0,1)$ subtracts $u_k$ from the lattice noise and $u_f/\sqrt{12}$ from the site noise modulo one. A double-precision finite-difference check agreed with the two JVP outputs to $4.0\times10^{-11}$.
Each fit used width 64, AdamW at $10^{-3}$, clipping at norm 10, 64 positive examples and 16 composition blocks per update. One, three and ten passes correspond to 1,696, 5,088 and 16,960 updates and 108,544, 325,632 and 1,085,440 positive examples per fit. Each drift update evaluated 64 independent negative and 64 live query structures. We chose one pass count per objective by the lowest mean validation geometry score at five draws on 128 fixed material IDs across three seeds. Flow selection used 16 Euler steps. The one-Euler diagnostic uses those same selected weights. The official test set had already been examined in the earlier study.
\begin{table}[!htbp]\centering\small
\caption{Validation-only checkpoint selection. The energy score measures forecast error in standardized log volume and three nearest-neighbor summaries. CPU hours sum process time over three complete ten-pass fits. Each fit used one CPU thread and zero GPU hours.}\label{t:exp-training-fits}
\begin{tabular}{lrrrrrr}\hline Objective & Parameters & Pass 1 & Pass 3 & Pass 10 & Selected & CPU h\\\hline
CrytalDrift & 24,377 & 2.790 & 2.822 & 2.857 & 1 & 0.55 \\
IID-positive drift & 24,377 & 2.795 & 2.822 & 2.854 & 1 & 0.55 \\
Compact flow & 24,377 & 2.737 & 2.761 & 2.717 & 10 & 0.20 \\
Compact MeanFlow & 24,383 & 2.860 & 2.919 & 2.845 & 10 & 0.62 \\
\hline\end{tabular}\end{table}
\begin{figure}[!htbp]\centering\includegraphics[width=.88\linewidth]{exp-training-curves.pdf}
\caption{Validation geometry score across the fixed training budgets. Thin lines show individual fit seeds and bold lines their arithmetic means. All arms use the same target IDs and five-draw input noises.}\label{fig:exp-training-curves}\end{figure}
Both drift samplers selected the first pass. At that budget the IID-minus-shuffled validation score difference was 0.0047 (paired 95\% interval $[-0.0039,0.0136]$). Training for ten passes increased both drift validation scores. We retained the original width for all four objectives. The two velocity objectives selected ten passes.
\begin{table}[!htbp]\centering\small\setlength{\tabcolsep}{3pt}
\caption{Composition-conditioned prediction on all 9,046 official MP-20 test targets, with matched cost measurements. Each raw draw is compared with its own target. RMS averages the lowest accepted error among recovered targets. Compact and retrieval accuracy rows equally weight three seeds. Each released model contributes one checkpoint. The warm batch-one median and fixed 10,000-request wall time use a dedicated GPU timing stream on the same 128-target prefix as Figure~\ref{fig:exp-gpu}. Retrieval runs on the host.}\label{t:exp-conditioned}
\resizebox{\linewidth}{!}{\begin{tabular}{lrrrrrrrrrr}\hline & & & \multicolumn{3}{c}{MR (\%)} & \multicolumn{3}{c}{Normalized RMS} & \multicolumn{2}{c}{Measured cost}\\
Method & $S$ & $kS$ (1/5/20) & $k{=}1$ & 5 & 20 & $k{=}1$ & 5 & 20 & B1 (ms) & min/10,000\\\hline
CrytalDrift & 1 & 1/5/20 & 6.79 & 20.33 & 32.91 & 0.382 & 0.371 & 0.346 & 15.01 & 0.119 \\
IID-positive drift & 1 & 1/5/20 & 6.85 & 20.59 & 32.98 & 0.383 & 0.371 & 0.346 & 15.05 & 0.088 \\
Compact MeanFlow & 1 & 1/5/20 & 8.76 & 21.41 & 31.86 & 0.388 & 0.362 & 0.339 & 13.70 & 0.102 \\
Flow, 1 Euler step & 1 & 1/5/20 & 7.49 & 18.00 & 25.04 & 0.403 & 0.389 & 0.359 & 12.68 & 0.059 \\
Flow, 16 Euler steps & 16 & 16/80/320 & 10.63 & 23.71 & 36.54 & 0.363 & 0.334 & 0.319 & 112.87 & 0.888 \\
Retrieval & 0 & 0/0/0 & 61.86 & 64.60 & 65.19 & 0.040 & 0.041 & 0.042 & 0.22 & 0.015 \\
Released uFlowCSP & 1 & 1/5/20 & 46.32 & 68.05 & 79.32 & 0.089 & 0.088 & 0.085 & 80.08 & 1.059 \\
CrystalFlow, 5 steps & 5 & 5/25/100 & 58.14 & 77.45 & 85.87 & 0.232 & 0.186 & 0.156 & 143.08 & 0.884 \\
\hline\end{tabular}}\end{table}
At twenty draws, differences in recovery (CrytalDrift minus the named method, percentage points with 95\% paired target-bootstrap intervals) were IID-positive drift -0.07 ($[-0.25,0.12]$), Compact MeanFlow +1.05 ($[0.64,1.46]$), Flow, 1 Euler step +7.87 ($[7.36,8.39]$), Flow, 16 Euler steps -3.63 ($[-4.06,-3.19]$).
\begin{table}[!htbp]\centering\small
\caption{Separating lattice and site errors. Shape is the Frobenius norm of the traceless log-stretch difference in the deposited basis. Volume is absolute log volume-per-atom error. Both average 20 draws per test target. Site RMS is a displacement on the unit fractional torus. For the three neighbor columns, we placed each generated motif in its reference lattice and measured mean absolute errors in the minimum, median and 90th percentile site-neighbor distances, in \AA, for five draws on the 128-target prefix.}\label{t:exp-shape}
\begin{tabular}{lrrrrrr}\hline Method & Shape & Volume & Site RMS & Minimum & Median & 90th pct.\\\hline
CrytalDrift & 0.728 & 0.189 & 0.290 & 1.359 & 0.856 & 0.717 \\
IID-positive drift & 0.727 & 0.188 & 0.289 & 1.357 & 0.856 & 0.715 \\
Compact MeanFlow & 0.676 & 0.091 & 0.290 & 1.358 & 0.861 & 0.716 \\
Flow, 1 Euler step & 0.570 & 0.089 & 0.307 & 1.460 & 0.961 & 0.782 \\
Flow, 16 Euler steps & 0.739 & 0.098 & 0.280 & 1.309 & 0.819 & 0.701 \\
\hline\end{tabular}\end{table}
The site RMS column gives a feasible species-matched RMS displacement on the unit fractional torus in the deposited basis. We initialize the translation by aligning a fixed atom of the rarest target species to each eligible candidate atom, alternate exact within-species assignments and mean-residual translation updates up to four times, and retain the smallest value. This gives an upper bound on the optimum over origins. It is evaluated on the same 128 targets and five draws as the site-neighbor columns.
The compact MeanFlow and drifting models had similar fractional-site displacement (0.290 and 0.290), while MeanFlow halved the mean absolute log-volume error (0.091 versus 0.189). The 16-step flow had the smallest fractional-site error (0.280) among the compact methods. Cell-shape error and site-placement error therefore favor different controls.
